%----------------------------------------------------------------------------------------
% Preambulo y Configuración
%----------------------------------------------------------------------------------------

\documentclass[
    11pt,
    spanish,
    singlespacing,
    parskip,
    headsepline,
    bookmarks=true,
    unicode=true,
    pdftoolbar=true,
    pdfmenubar=true,
    pdffitwindow=false,
    colorlinks=true,
    linkcolor=blue,
    citecolor=blue,
    urlcolor=blue
]{MastersDoctoralThesis}

\usepackage[utf8]{inputenc} % Codificación de entrada UTF-8
\usepackage[T1]{fontenc}    % Codificación de salida para caracteres especiales
\usepackage{graphicx}       % Manejo de gráficos
\usepackage{eso-pic}        % Permite agregar fondos
\usepackage{hyperref}       % Manejo de hipervínculos y marcadores

% Redefinición de caracteres problemáticos en marcadores
\hypersetup{
    pdftitle={Título del Documento},
    pdfauthor={Autor del Documento},
    pdfkeywords={Sistemas Embebidos, Internet de las Cosas, Inteligencia Artificial},
    pdfstartview={FitH},
    unicode=true,
    colorlinks=true,
    linkcolor=blue,
    citecolor=blue,
    urlcolor=blue
}

\pdfstringdefDisableCommands{%
  \def\texttt#1{#1}%
  \def\textbf#1{#1}%
  \def\textit#1{#1}%
  \def\"{\"}%
  \def\~{~}%
  \def\'{'}%
  \def\^{}%
  \def\textunderscore{\_} % Manejo del subrayado en marcadores
}


% Definir comandos requeridos por la clase
\newcommand{\degreename}{Maestría en Ciencias} % Cambia según tu título
\newcommand{\univname}{Universidad Nacional de Ejemplo} % Cambia según tu universidad
\newcommand{\keywordnames}{Palabras clave:}
%----------------------------------------------------------------------------------------
% Documento Principal
%----------------------------------------------------------------------------------------

\begin{document}

% Configuración de la portada
\posgrado{Carrera de Especialización en Sistemas Embebidos}
\keywords{Sistemas Embebidos, Internet de las Cosas, Bluethooth Mesh}

% Incluir la portada desde un archivo separado
\include{portada}

% Configuración del contenido preliminar
\frontmatter % Usar numeración romana para las páginas preliminares
\pagestyle{plain} % Estilo de encabezado simple

%----------------------------------------------------------------------------------------
% Resumen
%----------------------------------------------------------------------------------------

\begin{abstract}
\addchaptertocentry{\abstractname} % Agregar resumen al índice
El proyecto consiste en el diseño e implementación de un sistema Bluetooth Mesh para la gestión de sensores en invernaderos de Wentux Tecnoagro, con el fin de mejorar la escalabilidad y eficiencia de sus soluciones actuales.

Para su desarrollo fueron fundamentales los conocimientos adquiridos en la carrera, en áreas como programación de microcontroladores, sistemas operativos en tiempo real, ingeniería y testing de software.							
\end{abstract}

%----------------------------------------------------------------------------------------
% Índice
%----------------------------------------------------------------------------------------

\tableofcontents
\listoffigures
\listoftables

%----------------------------------------------------------------------------------------
% Capítulos
%----------------------------------------------------------------------------------------

\mainmatter % Iniciar numeración numérica para el contenido principal
\pagestyle{thesis} % Estilo de encabezado de tesis

% Incluir capítulos desde archivos separados
% Chapter 1

\chapter{Introducción general} % Main chapter title

\label{Chapter1} % For referencing the chapter elsewhere, use \ref{Chapter1} 
\label{IntroGeneral}

Todos los capítulos deben comenzar con un breve párrafo introductorio que indique cuál es el contenido que se encontrará al leerlo.  La redacción sobre el contenido de la memoria debe hacerse en presente y todo lo referido al proyecto en pasado, siempre de modo impersonal.

%----------------------------------------------------------------------------------------

% Define some commands to keep the formatting separated from the content 
\newcommand{\keyword}[1]{\textbf{#1}}
\newcommand{\tabhead}[1]{\textbf{#1}}
\newcommand{\code}[1]{\texttt{#1}}
\newcommand{\file}[1]{\texttt{\bfseries#1}}
\newcommand{\option}[1]{\texttt{\itshape#1}}
\newcommand{\grados}{$^{\circ}$}

%----------------------------------------------------------------------------------------
%\section{Introducción}

%----------------------------------------------------------------------------------------
\section{Estado del arte}

Revisión de tecnologías actuales para redes de sensores en invernaderos (cableadas e inalámbricas), comparando alternativas como Zigbee, LoRa, Wi-Fi y Bluetooth Mesh, con sus ventajas y limitaciones.



%----------------------------------------------------------------------------------------

\section{Motivación}

Explicación del problema (limitaciones del sistema cableado de Wentux Tecnoagro), la necesidad de escalabilidad y flexibilidad, y la relevancia de implementar una solución embebida basada en IoT.

%----------------------------------------------------------------------------------------

\section{Alcance y objetivos}

Definición de lo que cubre el trabajo y lo que queda fuera del alcance. Objetivo general y objetivos específicos.
\chapter{Introducción específica} % Main chapter title

\label{Chapter2}

%----------------------------------------------------------------------------------------
%	SECTION 1
%----------------------------------------------------------------------------------------
Todos los capítulos deben comenzar con un breve párrafo introductorio que indique cuál es el contenido que se encontrará al leerlo.  La redacción sobre el contenido de la memoria debe hacerse en presente y todo lo referido al proyecto en pasado, siempre de modo impersonal.

\section{Estructura del sistema}
\label{sec:ejemplo}

Descripción general del sistema propuesto: nodo central, nodos sensores/actuadores, topología Mesh, y cómo interactúan entre sí.


\section{Protocolos de comunicación}
Explicación de los protocolos usados: Bluetooth Mesh, UART y Wi-Fi/HTTP

\section{Requerimientos}
Requerimientos principales: funcionales, no funcionales y de validación
\chapter{Diseño e implementación} % Main chapter title

\label{Chapter3} % Change X to a consecutive number; for referencing this chapter elsewhere, use \ref{ChapterX}

Todos los capítulos deben comenzar con un breve párrafo introductorio que indique cuál es el contenido que se encontrará al leerlo.  La redacción sobre el contenido de la memoria debe hacerse en presente y todo lo referido al proyecto en pasado, siempre de modo impersonal.


%----------------------------------------------------------------------------------------
%	SECTION 1
%----------------------------------------------------------------------------------------
\section{Arquitectura de hardware}
Descripción del microcontrolador ESP32 y motivo por el que se seleccionó para este trabajo
\section{Arquitectura de firmware}
Descripción del software desarrollado y módulos principales. Estrategia para garantizar compatibilidad, escalabilidad y mantenibilidad.
\section{Arquitectura de interfaz de usuario}
Descripción de la interfaz de usuario, API implementada, y funcionalidades principales. Estrategia para garantizar compatibilidad, escalabilidad y mantenibilidad
\section{Documentación}
Descripción de la documentación entregada al cliente y obetivo: Documentación del código, guia del desarrollador, guia de usuario y pruebas del sistema




% Chapter Template

\chapter{Ensayos y resultados} % Main chapter title

\label{Chapter4} % Change X to a consecutive number; for referencing this chapter elsewhere, use \ref{ChapterX}
Todos los capítulos deben comenzar con un breve párrafo introductorio que indique cuál es el contenido que se encontrará al leerlo.  La redacción sobre el contenido de la memoria debe hacerse en presente y todo lo referido al proyecto en pasado, siempre de modo impersonal.

%----------------------------------------------------------------------------------------
%	SECTION 1
%----------------------------------------------------------------------------------------

\section{Banco de pruebas}
Entorno utilizado para la validación: número de nodos, configuraciones, herramientas de prueba, condiciones de simulación.

\section{Pruebas unitarias}
Ejemplos de pruebas unitarias realizadas sobre los módulos principales del firmware

\section{Pruebas funcionales}
\label{sec:pruebasHW}

Pruebas de comunicación entre nodos
Pruebas de latencia y tiempo de respuesta
Pruebas de desconexión y reconexión de nodos
Pruebas de interfaz web

% Chapter Template

\chapter{Conclusiones} % Main chapter title

\label{Chapter5} % Change X to a consecutive number; for referencing this chapter elsewhere, use \ref{ChapterX}
Todos los capítulos deben comenzar con un breve párrafo introductorio que indique cuál es el contenido que se encontrará al leerlo.  La redacción sobre el contenido de la memoria debe hacerse en presente y todo lo referido al proyecto en pasado, siempre de modo impersonal.


%----------------------------------------------------------------------------------------

%----------------------------------------------------------------------------------------
%	SECTION 1
%----------------------------------------------------------------------------------------

\section{Resultados obtenidos}
Resumen de logros alcanzados en comparación con los objetivos iniciales. Beneficios obtenidos para Wentux Tecnoagro y análisis de limitaciones.

\section{Trabajo futuro}

Integración con dispositivos reales de la empresa, pruebas en invernadero, seguridad avanzada, soporte en la nube.


%----------------------------------------------------------------------------------------
% Apéndices
%----------------------------------------------------------------------------------------

\appendix

% Incluir apéndices desde archivos separados si es necesario
%\include{Appendices/AppendixA}

%----------------------------------------------------------------------------------------
% Bibliografía
%----------------------------------------------------------------------------------------

\renewcommand{\bibname}{Bibliografía} % Para asegurarte de que el título sea correcto
\phantomsection % Necesario para que el enlace del marcador sea correcto

\printbibliography[heading=bibintoc]

\end{document}






